\documentclass{beamer}
\usetheme{Madrid} % Tema clássico e limpo (outras opções: Metropolis, Frankfurt)
\usecolortheme{default}
\usepackage{booktabs}
\usepackage{graphicx}
% Informações da Capa
\title[ARC-AGI-1 e CEGIS]{Análise do Impacto de Variações de Prompting e Feedback Semântico no ARC-AGI-1}
\author[Cieslak, Machado]{Fabio Luiz da C. Cieslak \and Matheus Machado}
\institute[UFRGS]{Universidade Federal do Rio Grande do Sul}
\date{Setembro 2026}

\begin{document}

% 1º Slide: Capa automática
\begin{frame}
    \titlepage
\end{frame}

% 2º Slide: Exemplo de conteúdo com tópicos
\begin{frame}{Introdução ao Problema}
    \begin{itemize}
        \item O benchmark ARC-AGI-1 avalia raciocínio indutivo e generalização.
        \item Modelos costumam recorrer a sobreajustes espúrios (ramificações \textit{if/else}).
        \item \textbf{Nossa abordagem:} Utilizar a síntese indutiva guiada por contraexemplos (CEGIS) com regras Anti-Trapaça.
    \end{itemize}
\end{frame}
% 3º Slide: Metodologia
\begin{frame}{Metodologia}
    \begin{itemize}
        \item \textbf{Avaliação:} 100 tarefas do conjunto de treino do ARC-AGI-1.
        \item \textbf{Modelos avaliados:} 9 modelos de linguagem (fronteira, intermediários e compactos).
        \item \textbf{Métricas Analisadas:}
        \begin{itemize}
            \item Acurácia (sucesso estrito no teste cego).
            \item Custo iterativo (requisições por tarefa).
            \item Taxonomia de erros (Recuperação vs. Regressão).
        \end{itemize}
        \item \textbf{Inferência Estatística:} Reamostragem Bootstrap (10.000 iterações).
    \end{itemize}
\end{frame}

% 4º Slide: Resultados (Tabela 1)
\begin{frame}{Resultados: Desempenho e Eficiência}
    \begin{table}[ht]
        \centering
        \resizebox{0.95\textwidth}{!}{%
        \begin{tabular}{l ccccc}
            \toprule
            Modelo & CEGIS & Anti-Trapaça & $\Delta$ & Req./Tarefa & Recuperação \\
            \midrule
            \texttt{Grok 4.6} & 92 & 92 & +0 & 1.13 & 4 \\
            \texttt{DeepSeek V4 Flash 0731} & 86 & 86 & +0 & 1.25 & 12 \\
            \texttt{GPT-5.6 Luna} & 77 & 79 & +2 & 2.00 & 19 \\
            \texttt{Claude Sonnet 5} & 62 & 60 & -2 & 2.90 & 30 \\
            \texttt{Gemini 3.5 Flash Lite} & 43 & 45 & +2 & 3.23 & 22 \\
            \texttt{Gemini 3.1 Flash Lite} & 31 & 24 & -7 & 4.07 & 19 \\
            \texttt{GLM 5.3 Flash} & 19 & 17 & -2 & 4.42 & 9 \\
            \texttt{Leanstral 1.5.1} & 7 & 9 & +2 & 4.69 & 4 \\
            \texttt{GPT-5.4 Nano} & 5 & 3 & -2 & 4.89 & 4 \\
            \bottomrule
        \end{tabular}%
        }
    \end{table}
    
    \vspace{0.3cm}
    \small \textbf{Destaque:} A regra Anti-Trapaça alavancou os modelos intermediários. Foram observadas \textbf{zero regressões} em todos os testes.
\end{frame}
% 5º Slide: Validação Estatística (Tabela 2)
\begin{frame}{Resultados: Validação Estatística e Erros}
    \begin{table}[ht]
        \centering
        \resizebox{0.95\textwidth}{!}{%
        \begin{tabular}{lcccc}
            \toprule
            Modelo & $p_{\text{boot}}$ & IC 95\% Bootstrap & Sobreajuste Espúrio & Teto Repres. \\
            \midrule
            \texttt{Grok 4.6} & 0.0431 & [1.0\%, 8.0\%] & 8 & 0 \\
            \texttt{DeepSeek V4 Flash 0731} & 0.0007 & [6.0\%, 19.0\%] & 10 & 4 \\
            \texttt{GPT-5.6 Luna} & 0.0001 & [12.0\%, 27.0\%] & 10 & 13 \\
            \texttt{Claude Sonnet 5} & 0.0001 & [21.0\%, 39.0\%] & 7 & 31 \\
            \texttt{Gemini 3.5 Flash Lite} & 0.0001 & [14.0\%, 30.0\%] & 10 & 47 \\
            \texttt{Gemini 3.1 Flash Lite} & 0.0001 & [12.0\%, 27.0\%] & 7 & 62 \\
            \texttt{GLM 5.3 Flash} & 0.0052 & [4.0\%, 15.0\%] & 0 & 81 \\
            \texttt{Leanstral 1.5.1} & 0.0457 & [1.0\%, 8.0\%] & 0 & 93 \\
            \texttt{GPT-5.4 Nano} & 0.0496 & [1.0\%, 8.0\%] & 0 & 95 \\
            \bottomrule
        \end{tabular}%
        }
    \end{table}
    
    \vspace{0.3cm}
    \small \textbf{Análise:} Ganhos estatisticamente significantes ($p < 0.05$) em toda a escala de modelos. Modelos compactos esbarram no \textbf{Teto de Representação}.
\end{frame}

% 6º Slide: Conclusão
\begin{frame}{Conclusões}
    \begin{itemize}
        \item \textbf{Eficácia do CEGIS:} A retroalimentação semântica baseada em contraexemplos gera ganhos consistentes (com avanços de até +30\%).
        \item \textbf{Segurança Estrutural:} Preservação total de acertos prévios ao longo dos ciclos iterativos (estrita monotonicidade).
        \item \textbf{Impacto do Anti-Trapaça:} Mitigou o sobreajuste espúrio, trazendo benefícios visíveis para modelos intermediários.
        \item \textbf{Rigor Metodológico:} O Bootstrap permitiu discernir aprimoramentos reais do ruído estocástico, criando um modelo de validação escalável.
    \end{itemize}
\end{frame}

% 7º Slide: Encerramento
\begin{frame}
    \begin{center}
        \Huge \textbf{Obrigado} \\
        \vspace{0.8cm}
        \normalsize Código e Dados: \texttt{github.com/Jubarte27/arc-agi-1}
    \end{center}
\end{frame}
\end{document}